\section{六次元および七次元の幾何}\label{sec:pg-six-seven}
\subsection{可能な階数}
定理\ref{thm:pg-normal}により次の表を得る。
\begin{center}
\begin{tabular}{ccccc}\toprule
$n$ & $k$ & 一般階数 & $m$ & Poisson 超平面\\\midrule
6&1&2&5&$\PP^4$\\
6&2&4&2&$\PP^1$\\
7&1&2&6&$\PP^5$\\
7&2&4&3&$\PP^2$\\
7&3&6&0&なし\\\bottomrule
\end{tabular}
\end{center}
全双ベクトル場の空間の次元は $\binom83=56$、$\binom93=84$。
階数二の構造の射影化は $\operatorname{Gr}(3,8)$、$\operatorname{Gr}(3,9)$ で、次元は十五、十八である。

\subsection{階数二の零点：一般型および Calabi--Yau 完全交叉}
$k=1$ とし、active 座標を $x,y,z$ と書く。$m=n-1$ である。
$C=X\cap\{u=0\}$、$\Gamma=V(F_x,F_y,F_z)$、$Y=\Gamma\cap X$ とする。
\begin{proposition}\label{prop:pg-rank2}
固定した任意の滑らかな $X$ に対し、一般の三次元部分空間を active ブロックに選ぶと、
$C,\Gamma,Y$ は全て滑らかで、スキームとして
$D_0(\pi)=C\sqcup\Gamma$、$D_0(\pi_X)=C\sqcup Y$ である。
さらに
\begin{equation}
\Gamma:(2,2,2),\quad\deg\Gamma=8,\quad K_\Gamma=(4-n)H,
\qquad
Y:(2,2,2,3),\quad\deg Y=24,\quad K_Y=(7-n)H.
\end{equation}
\end{proposition}
\begin{proof}
全偏微分の張る二次式の線形系は射影空間上基点をもたない。
Bertini を三回適用すると一般の三線形結合の零点 $\Gamma$ は滑らかとなる。
同じ線形系を滑らかな $X$ に制限すれば $Y$ も滑らかとなる。
一般の平面切断 $C$ の滑らかさも Bertini による。
これらは active 部分空間の Grassmann 多様体上の非空開条件なので同時に満たせる。
命題\ref{prop:pg-deg-support}、\ref{prop:pg-zero-ci}と随伴公式を適用する。
\end{proof}
$n=6$ では $\Gamma$ は $-K_\Gamma=2H$ の次数八の Fano 四次元多様体、
$Y$ は $K_Y=H$、$K_Y^3=24$ の一般型三次元多様体となる。
$n=7$ では $\Gamma$ は $-K_\Gamma=3H$ の Fano 五次元多様体、
$Y$ は次数二十四の Calabi--Yau 四次元多様体である。
ここで Calabi--Yau とは $K_Y\simeq\OO_Y$ および $H^i(Y,\OO_Y)=0$（$0<i<4$）を意味し、
Lefschetz 定理から単連結性も従う。
$h^{p,q}=\dim H^q(\Omega^p)$、$p_g=h^0(\omega)$、$e$ を位相的 Euler 数と書くと
\begin{equation}\begin{aligned}
n=6:&\quad p_g(Y)=8,\quad h^{1,1}(Y)=1,\quad h^{2,1}(Y)=246,\quad e(Y)=-504,\\
n=7:&\quad h^{1,1}(Y)=1,\quad h^{2,1}(Y)=0,\quad h^{3,1}(Y)=183,\quad h^{2,2}(Y)=780,\quad e(Y)=1152.
\end{aligned}\end{equation}
計算は\S\ref{sec:pg-numerical}で与える。

\subsection{階数四と中心変数}
$F=f(a,u)+g(b,u)$、$a,b\in\CC^3$ とし、
$A_i=f_{a_i}$、$B_j=g_{b_j}$、$\rho_f=\sum_i a_iA_i$、$\rho_g=\sum_jb_jB_j$ とおく。
命題\ref{prop:pg-deg-ideal}は
\begin{equation}\begin{aligned}
I_{D_2(\pi)}&=(u_\alpha A_iB_j,\rho_gA_i,\rho_fB_j),\\
I_{D_0(\pi)}&=(u_\alpha A_i,u_\alpha B_j,\rho_f,\rho_g,b_jA_i,a_iB_j)
\end{aligned}\end{equation}
を与える。
$f_0,g_0$ が滑らかな平面三次曲線 $E_f,E_g$ を定めるとし、
$\Gamma_f=V(A_1,A_2,A_3)$、$\Gamma_g=V(B_1,B_2,B_3)$、
$J=V(u,f_0,g_0)$ とおく。
$J$ は二曲線の三次元 join であり、
\begin{equation}
|D_2(\pi)|=|\Gamma_f|\cup|\Gamma_g|\cup J,
\qquad |D_2(\pi_X)|=|\Sigma_f|\cup|\Sigma_g|\cup J,
\quad\Sigma_f=\Gamma_f\cap X,
\end{equation}
および同様の $\Sigma_g$ を得る。
\begin{proposition}\label{prop:pg-cubic-surface-family}
$\Sigma_f,\Sigma_g$ は型 $(2,2,2,3)$、次元 $m+1$、次数二十四の完全交叉であり、
一般のデータでは既約である。
$K_f=V(A_1,A_2,A_3)\subset\PP(a,u)$ および同様の $K_g$ を底として、
それぞれの双有理モデルは平坦な三次曲面族をなす。
相対 Jacobian イデアルで定まる fiber の特異スキームは $Z=\Delta\cap X$ である。
一般 fiber は滑らかだが全 fiber が滑らかではない。
\end{proposition}
\begin{proof}
$\Gamma_f$ は $K_f$ と反対側の $\PP_b^2$ の join であり、その斉次環は
$\CC[a,u]/(A)$ に変数 $b$ を加えた環である。
$F$ の最高 $b$ 次数部分 $g_0(b)$ は非零で係数は $\CC$ に属する。
従って随伴次数環での $g_0$ の乗法は単射であり、元の環で $F$ は非零因子となる。
補題\ref{lem:pg-critical-cover}から完全交叉性と数値を得る。

一般のデータを取り $K_f$ を滑らかとする。
join の頂点の解消として、直線をパラメータ化する
$\mathbb P_f=\PP_{K_f}(\OO_{K_f}(-1)\oplus\OO_{K_f}^{\oplus3})\to\Gamma_f$ を取る。
底の代表 $(a,u)$ を使う表示では $[t:b]\mapsto[ta:tu:b]$ であり、strict transform は
\begin{equation}
t^3f(a,u)+g(b,tu)=0.\label{eq:pg-surface-fiber}
\end{equation}
代表変更に伴い $t$ が逆に変わるので式は貼り合わさる。
$t=0$ に全成分は含まれず、余分な例外成分はない。
各 fiber の式は $t=0$ で非零の $g_0(b)$ となり、常に三次曲面を定める。
相対超曲面の完全列は fiber への制限後も左端が単射なので、族は平坦である。
$t=0$ では $g_0=0$ が滑らかなので相対特異点はない。
$t=1$ では相対特異イデアルは $(F,B_1,B_2,B_3)$ で、底の $(A_1,A_2,A_3)$ を合わせると正確に $Z$ となる。
$\Delta\to K_f$ は有限で $\dim Z=m-2$。
$K_f$ は一般に滑らかな連結完全交叉、従って既約で次元 $m-1$ なので $Z$ は底を支配しない。
一般 fiber は滑らかな三次曲面で幾何学的に既約。
平坦性により垂直な既約成分はなく、全空間とその像 $\Sigma_f$ は既約である。
$g$ 側も同様であり、$\deg Z=192>0$ から特異 fiber は存在する。
\end{proof}
一般のデータについての数値は次の通りである。
\begin{center}
\begin{tabular}{lll}\toprule
 & $n=6,m=2$ & $n=7,m=3$\\\midrule
$K_f,K_g$ & 次数8、種数5の曲線 & 次数8の K3 曲面\\
$\dim\Sigma_f$ & 3 & 4\\
$\Delta$ & 次数64、種数129の曲線 & 次数64、$K_\Delta=3H$ の曲面\\
$Z$ & 192個の被約点 & 次数192、種数577の滑らかな曲線\\
$D_0(\pi_X)$ & $E_f\sqcup E_g\sqcup Z$ & $E_f\sqcup E_g\sqcup Z$\\\bottomrule
\end{tabular}
\end{center}
$n=7$ の $\Delta$ では
$K_\Delta^2=576$、$c_2(\Delta)=768$、$\chi(\OO_\Delta)=112$、$p_g(\Delta)=111$、$h^1(\OO_\Delta)=0$。
両次元の一般葉の閉包は $\PP^6$ 内の型 $(3,3)$ の四次元完全交叉であり、
$X$ 上の底はそれぞれ $\PP(1,1,3)$、$\PP(1,1,1,3)$。
対数表示の反標準因子の次数は $(1,1,3,3)$、$(1,1,1,3,3)$ である。

\subsection{七次元の階数六}
ここでは $m=0$、$F=f(a)+g(b)+h(c)$、$a,b,c\in\CC^3$ である。
$X$ の滑らかさは各 $f,g,h$ が滑らかな平面三次曲線を定めることと同値である。
一方が特異なら対応する座標平面に $X$ の特異点が生じ、
逆に各 gradient の共通零点が原点だけなら全 gradient についても同じだからである。
曲線を $E_f,E_g,E_h$、二つずつの join を $J_{fg},J_{fh},J_{gh}$、
三曲線の join を $J_3=V(f,g,h)$ とする。
それぞれの次元は三、五である。
一次 Casimir と Poisson 超平面は存在しないが、$f,g,h$ の線形結合は Casimir である。
最高外積は
\begin{equation}\begin{aligned}
\pi^3&\longleftrightarrow\iota_E(df\wedge dg\wedge dh)\\
&=3(f\,dg\wedge dh-g\,df\wedge dh+h\,df\wedge dg)\\
&=3fgh\,d\log(g/f)\wedge d\log(h/f)\in H^0(\PP^8,\Omega^2(9)).
\end{aligned}\end{equation}
$[f:g:h]:\PP^8\dashrightarrow\PP^2$ は $X$ 上で直線 $f+g+h=0$ に値を取る。
$f/g$ は有理 Casimir で、一般六次元葉の閉包は $\PP^8$ 内の型 $(3,3)$ の完全交叉。
$X$ 上の最高外積に対応する一形式は
$f\,dg-g\,df\in H^0(X,\Omega_X^1(6))$ である。
退化の台は
\begin{equation}\begin{aligned}
|D_4(\pi)|&=J_3\cup V(a)\cup V(b)\cup V(c),\\
|D_2(\pi)|&=J_{fg}\cup J_{fh}\cup J_{gh}\cup\PP_a^2\cup\PP_b^2\cup\PP_c^2,\\
D_0(\pi)&=E_f\sqcup E_g\sqcup E_h.
\end{aligned}\end{equation}
$X$ 上では第一行の三つの $\PP^5$ が滑らかな三次四次元多様体となり、
第二行の三つの平面はそれぞれ曲線に制限され、既に join に含まれる。
従って $|D_2(\pi_X)|=J_{fg}\cup J_{fh}\cup J_{gh}$。

\subsection{非被約構造と primary 分解}
$A_i=f_{a_i}$、$B_j=g_{b_j}$、$C_l=h_{c_l}$ とし、
$\mathfrak a=(A_1,A_2,A_3)$、$\mathfrak b=(B_1,B_2,B_3)$、$\mathfrak c=(C_1,C_2,C_3)$ と書く。
\begin{proposition}\label{prop:pg-D4-primary}
周囲では $I_{D_4(\pi)}=(fB_jC_l,gA_iC_l,hA_iB_j)$。
$X$ 上ではイデアル層として
\begin{equation}
I_{D_4(\pi_X)}=I_{J_3}\cap\mathfrak a\cap\mathfrak b\cap\mathfrak c.\label{eq:pg-D4-primary}
\end{equation}
$J_3$ は被約成分で、他の三成分は三次四次元多様体を台とする primary な厚化である。
その一般点での横断長は八である。
\end{proposition}
\begin{proof}
生成元は命題\ref{prop:pg-deg-ideal}の $j=3$ である。
punctured affine cone の $A_i\neq0$ の開集合で $f=-g-h$ とおけば
$I_{D_4}=(g\mathfrak c,h\mathfrak b)$。
この開集合から $b,c$ 空間への写像は $a_i$ 方向の偏微分が単元なので滑らか、従って平坦である。
独立変数環 $\CC[b,c]$ では
\begin{equation}
(g\mathfrak c,h\mathfrak b)=(g,h)\cap\mathfrak b\cap\mathfrak c.
\end{equation}
Euler 等式から $g\in\mathfrak b$、$h\in\mathfrak c$ なので一方の包含は明らか。
逆に $gr+hs$ が右辺に属せば、$\mathfrak b$ を法として非零多項式 $h(c)$ の乗法は単射だから $s\in\mathfrak b$。
同様に $r\in\mathfrak c$ となる。
平坦基底変換は有限交叉と可換し、全点ではいずれかの gradient 成分が単元となるので、これらを貼り合わせて\eqref{eq:pg-D4-primary}を得る。

例えば $\CC[b]/\mathfrak b$ は Hilbert 級数 $(1+t)^3$、長さ八の局所 Artin 環である。
$\mathfrak b$ による錐の商は
$\CC[a,c]/(f+h)\otimes_\CC\CC[b]/\mathfrak b$。
第一因子は滑らかな射影三次四次元多様体の錐の整域である。
第二因子の冪零根基の冪で filter すると、被約商で非零の元は非零因子であると分かる。
従って primary で、一般点の横断長は八。
$J_3$ は独立な三曲線の錐の積の射影化なので被約既約である。
\end{proof}
\begin{proposition}\label{prop:pg-D2-primary}
周囲では
\begin{equation}
I_{D_2(\pi)}=(fB_j,gA_i,fC_l,hA_i,gC_l,hB_j,
 c_lA_iB_j,b_jA_iC_l,a_iB_jC_l).
\end{equation}
$X$ 上では
\begin{equation}
I_{D_2(\pi_X)}=I_{J_{fg}}\cap I_{J_{fh}}\cap I_{J_{gh}}
\cap(\mathfrak a+\mathfrak b)\cap(\mathfrak a+\mathfrak c)\cap(\mathfrak b+\mathfrak c).
\label{eq:pg-D2-primary}
\end{equation}
最後の三項はそれぞれ $E_h,E_g,E_f$ を台とする埋め込まれた primary 成分であり、各厚化の横断長は六十四である。
\end{proposition}
\begin{proof}
生成元は $j=2$ の公式による。
$A_i\neq0$ 上で $F=0$ を使うと
$I_{D_2}=(g,h,c_lB_j,b_jC_l)$。
$\mathfrak m_b=(b_1,b_2,b_3)$、$\mathfrak m_c=(c_1,c_2,c_3)$ とおくと、独立変数環で
\begin{align}
(\mathfrak m_c,g)\cap(\mathfrak m_b,h)&=(g,h,\mathfrak m_b\mathfrak m_c),\label{eq:pg-D2-reduced}\\
(g,h,\mathfrak m_b\mathfrak m_c)\cap(\mathfrak b+\mathfrak c)
&=(g,h,\mathfrak m_c\mathfrak b+\mathfrak m_b\mathfrak c).\label{eq:pg-D2-intersection}
\end{align}
第一式は $(g,h)$ で割って各因子を定数部分と正次数部分に分解すれば分かる。
第二式は $b,c$ 二重次数で確認する。
両次数が正なら $\mathfrak b+\mathfrak c$ のその成分は
$\mathfrak m_c\mathfrak b+\mathfrak m_b\mathfrak c$ に等しく、一方の次数が零なら第一式のイデアルは $(g)$ または $(h)$ であり、Euler 等式により $\mathfrak b+\mathfrak c$ に入る。
平坦基底変換により局所的に
$I_{D_2}=(\mathfrak m_c,g)\cap(\mathfrak m_b,h)\cap(\mathfrak b+\mathfrak c)$ となる。
初めの二項はそこでの $J_{fg},J_{fh}$ のイデアルである。
他の gradient 開集合でも同様に計算し、錐から $X$ へ降ろせば\eqref{eq:pg-D2-primary}を得る。

$\mathfrak b+\mathfrak c$ の商は
$\CC[a]/(f)\otimes\CC[b]/\mathfrak b\otimes\CC[c]/\mathfrak c$。
前命題と同じ filtration により primary で、横断長は $8\cdot8=64$ となる。
この成分は冗長ではない。
$b_jc_l$ は\eqref{eq:pg-D2-reduced}の右辺に属すが二重次数 $(1,1)$ をもち、$\mathfrak b+\mathfrak c$ に入らない。
台 $E_f$ は二つの join に含まれるので埋め込み成分である。他も同様。
\end{proof}
\begin{remark}\label{rem:pg-length49}
六十四は明示した primary 厚化の横断長であり、スキームと被約化との差の長さではない。
$E_f$ の一般点では
$\ker(\OO_{D_2(\pi_X)}\to\OO_{D_2(\pi_X)_{\mathrm{red}}})$ の横断長は四十九である。
実際、$Q_g=\CC[b]/\mathfrak b$、$Q_h=\CC[c]/\mathfrak c$ の極大イデアルを
$\mathfrak n_g,\mathfrak n_h$ とすると、\eqref{eq:pg-D2-reduced}、\eqref{eq:pg-D2-intersection}により核は横断方向で $\mathfrak n_g\otimes\mathfrak n_h$ に等しい。
各 Jacobian 環は長さ八で定数部分が一次元なので、長さは $7^2=49$ となる。
\end{remark}

